\demo{绪论应包括：本研究课题的学术背景及其理论与实际意义；国内外文献综述；相关领域的研究进展及成果、存在的不足或有待深入研究的问题；本研究课题的来源及主要研究内容；本文结构说明。插图或插表之前，文中必须有关于本插图或表的提示。}

\subsection{研究背景}

毕业设计（论文）是本科教育的重要实践环节，其撰写质量直接关系到学生的综合能力培养。一篇格式规范、排版美观的论文不仅能够提升阅读体验，更能体现作者严谨的学术态度。目前，大多数学生使用 Microsoft Word 进行论文撰写，但在面对复杂的格式要求时，往往需要耗费大量精力进行手动调整，容易出现前后不一致的问题\cite{knuth1984texbook}。

\LaTeX{} 作为一种专业的排版系统，凭借其强大的宏包生态与稳定的输出质量，在学术出版领域得到了广泛应用。将 \LaTeX{} 应用于毕业设计论文的排版，可以有效解决 Word 排版中的诸多痛点，如公式编号、交叉引用、参考文献管理、自动生成目录等\cite{lamport1994latex}。

\subsection{国内外研究现状}

\subsubsection{国外研究现状}

国外高校普遍重视论文排版的标准化工作。许多知名大学（如 MIT、Stanford、Cambridge 等）均提供了官方认证的 \LaTeX{} 论文模板，覆盖本科、硕士及博士各个层次\cite{mittelbach2004latex}。这些模板通常由学校的研究生院或图书馆维护，并配有详尽的使用文档，极大地降低了学生的排版负担。

\subsubsection{国内研究现状}

国内高校近年来也逐渐开始推广 \LaTeX{} 论文模板。清华大学、北京大学、上海交通大学等均已推出各自的非官方或半官方 \LaTeX{} 模板，形成了活跃的开发者社区\cite{li2019thuthesis}。然而，对于地方本科院校而言，针对本校毕业设计模板定制的 \LaTeX{} 解决方案仍然较为匮乏。

\subsection{研究内容与论文结构}

本文的主要研究内容包括：

1.深入分析学校毕业设计 Word 模板的排版规范，提取各级标题、正文、图表、参考文献等元素的格式参数。

2.基于 ctex、titlesec、gbt7714 等宏包，开发可直接使用的文档类文件（.cls）。

3.通过实际论文撰写案例，验证模板的排版效果与稳定性。

本文的后续章节安排如下：第 2 章介绍相关技术与理论基础；第 3 章阐述模板系统的设计与实现细节；第 4 章展示实验结果与分析；最后总结全文工作并展望未来改进方向。
